\unit{Exposé 212 -- Préschémas quotients}
\sid{212}
\src{99}
\seminar{13e année, 1960/61, n\textsuperscript{o} 212\\\hfill Février 1961}
\paperhead{TECHNIQUES DE CONSTRUCTION ET THÉORÈMES D'EXISTENCE EN GÉOMÉTRIE ALGÉBRIQUE\\
III : PRÉSCHEMAS QUOTIENTS}%
  {par Alexander GROTHENDIECK}

\subsection*{\textnormal{\uline{Introduction}.}}
\addcontentsline{toc}{subsection}{Introduction}

Les problèmes traités dans le présent exposé diffèrent de ceux envisagés dans
les deux précédents, en ce qu'on essaye de représenter certains foncteurs
covariants, et non plus contravariants, de schémas variables. Le procédé de
passage au quotient est cependant essentiel dans beaucoup de questions de
construction en géométrie algébrique, y compris celles des exposés I et II
([1], [2]). Ainsi, la question de \uline{l'effectivité d'une donnée de descente}
sur un $\symup{T}$-préschéma $\symup{X}$, relativement à un morphisme fidèlement plat et
quasi-compact $\symup{T\to S}$, équivaut à la question de l'existence d'un
quotient de $\symup{X}$ (satisfaisant les propriétés raisonnables examinées plus bas),
pour la relation d'équivalence plate dans $\symup{X}$ définie par la donnée de descente ;
les questions soulevées dans [1], A, \S~2 c se résoudront sans doute en même
temps que les questions posées dans le paragraphe numéro 2 du présent exposé.
De même, le \uline{schéma de Picard} (pour la définition, cf. [2], C, \S~3)
d'un $\symup{S}$-schéma $\symup{X}$ peut se définir de diverses manières comme quotient de
certains autres schémas (de diviseurs positifs, ou d'immersions dans un
projectif) par des relations d'équivalence plates, la définition et la
construction de ces schémas auxiliaires étant d'ailleurs techniquement plus
simple : ce sont en effet des schémas du type
$\symup{\underline{\operatorname{Hom}}_S(X,Y)}$ et variantes définis dans [2], C,
\S~2, et dont la construction sera l'objet de l'exposé suivant (sous des
hypothèses de projectivité convenables). Combinant donc les résultats du
présent exposé et du suivant, on arrive à la construction des schémas de
Picard, sous des hypothèses convenables.

Le problème du passage au quotient dans les préschémas offre encore plusieurs
questions non résolues. La plus importante est mentionnée dans le paragraphe
8. Elle reste actuellement le seul obstacle à la construction des
\uline{schémas de modules sur les entiers pour les courbes de genre arbitraire},
les \uline{variétés abéliennes polarisées}, etc. C'est dire que sa solution mérite les
efforts des spécialistes des groupes algébriques.

\src{100}
\subsection*{\textnormal{1. \uline{Relations d'équivalence, relations d'équivalence effectives}.}}
\addcontentsline{toc}{subsection}{1. Relations d'équivalence, relations d'équivalence effectives}

Soient $\symup{\text{\uwave{C}}}$ une catégorie, et $\symup{X}$ un objet de $\symup{\text{\uwave{C}}}$. Un couple
de morphismes
\[
  \symup{p_1,p_2:R\rightrightarrows X}
\]
est dit \fquote{\uline{couple d'équivalence}} dans $\symup{\text{\uwave{C}}}$, de \uline{but} $\symup{X}$ et
de \uline{source} $\symup{R}$, si pour tout objet $\symup{T}$ de $\symup{\text{\uwave{C}}}$, les deux
applications correspondantes
\[
  \symup{p_1(T),p_2(T):R(T)\rightrightarrows X(T)}
\]
(où pour tout objet $\symup{Y}$ de $\symup{\text{\uwave{C}}}$, on pose
$\symup{Y(T)=\operatorname{Hom}(T,Y)}$) définissent une application
\[
  \symup{R(T)\to X(T)\times X(T)}
\]
induisant une bijection de $\symup{R(T)}$ sur le graphe d'une relation d'équivalence
dans l'ensemble $\symup{X(T)}$. On introduit entre les couples d'équivalences de but
$\symup{X}$ une relation d'équivalence évidente, une classe d'équivalence pour cette
dernière est appelée une $\symup{\text{\uwave{C}}}$\uline{-relation d'équivalence} dans $\symup{X}$,
ou simplement une relation d'équivalence si aucune confusion n'en résulte.

Si $\symup{X\times X}$ existe, la donnée d'une relation d'équivalence dans $\symup{X}$
équivaut à la donnée d'un sous-objet $\symup{R}$ de $\symup{X\times X}$, tel que, pour tout
objet $\symup{T}$ de $\symup{\text{\uwave{C}}}$, le sous-ensemble de $\symup{(X\times X)(T)=X(T)\times X(T)}$ qui correspond à $\symup{R(T)}$ soit le graphe d'une relation d'équivalence dans
$\symup{X(T)}$. Désignant par $\symup{p_1}$ et $\symup{p_2}$ les morphismes de $\symup{R}$ dans $\symup{X}$ induits
par les projections $\symup{\operatorname{pr}_1}$ et $\symup{\operatorname{pr}_2}$, la
condition précédente exprime que $\symup{(p_1,p_2)}$ est un couple d'équivalence. On
peut aussi exprimer diagrammatiquement dans $\symup{\text{\uwave{C}}}$ (sous réserve de
l'existence de $\symup{X\times X}$ et du produit fibré
$\symup{(R,p_2)\times_X(R,p_1)}$) les axiomes d'une relation d'équivalence au sens
ensembliste pour les $\symup{R(T)}$ dans les $\symup{X(T)}$, conformément au principe général
expliqué dans [2], A, \S~1. Nous n'en aurons pas besoin.

Chaque fois qu'on a un couple de morphismes $\symup{(p_1,p_2)}$ de même source $\symup{R}$ et
de même but $\symup{X}$, on peut définir le \uline{conoyau} du couple comme l'objet $\symup{Y}$
de $\symup{\text{\uwave{C}}}$ qui représente le foncteur covariant en $\symup{Z}$ :
\[
  \symup{\operatorname{Hom}_{p_1,p_2}(X,Z),}
\]
ensemble des morphismes $\symup{u}$ de $\symup{X}$ dans $\symup{Z}$ tels que $\symup{up_1=up_2}$. Si $\symup{Y}$
existe, il est déterminé à un isomorphisme unique près. On le notera
$\symup{Y/(p_1,p_2)}$ ou, par abus de notation, $\symup{Y/R}$, cette dernière notation étant
surtout employée lorsque
\src{101}$\symup{(p_1,p_2)}$ est un couple d'équivalence : il est d'usage alors
d'identifier dans les notations la relation d'équivalence définie par le
couple, et $\symup{R}$. Noter que si on considère $\symup{Y}$ comme un quotient de $\symup{X}$, il ne
dépend en effet que de la relation d'équivalence définie par le couple
$\symup{(p_1,p_2)}$.

Partons maintenant d'un morphisme
\[
  \symup{f:X\to Y}
\]
qui permet donc de considérer $\symup{X}$ comme un \fquote{objet au-dessus de $\symup{Y}$}, et
supposons que le produit fibré
\[
  \symup{\text{\uwave{R}}(f)=X\times_Y X}
\]
existe. Soient $\symup{p_1}$ et $\symup{p_2}$ ses deux projections. Alors $\symup{(p_1,p_2)}$ est un
couple d'équivalence, dit \uline{associé} au morphisme $\symup{f}$. Il définit donc une
relation d'équivalence, dite associée à $\symup{f}$.

On dit qu'un couple de morphismes $\symup{(p_1;p_2)}$ de but $\symup{X}$, de même source $\symup{R}$,
est un \uline{couple d'équivalence effectif}, si
\begin{enumerate}
  \item[(i)] le conoyau $\symup{X/(p_1,p_2)=Y}$ existe
  \item[(ii)] le produit fibré $\symup{X\times_Y X}$ existe
  \item[(iii)] le morphisme $\symup{R\to X\times_Y X}$ de composantes $\symup{p_1}$ et
  $\symup{p_2}$ est un isomorphisme.
\end{enumerate}
Alors le couple $\symup{(p_1,p_2)}$ est bien un couple d'équivalence. On dit aussi que
la relation d'équivalence qu'il définit est une \uline{relation d'équivalence
effective}.

On dit qu'un morphisme $\symup{f:X\to Y}$ est un \uline{épimorphisme
effectif} si
\begin{enumerate}
  \item[(i)] le produit fibré $\symup{X\times_Y X=R}$ existe;
  \item[(ii)] le quotient $\symup{X/(p_1,p_2)}$ existe, où $\symup{p_1}$ et $\symup{p_2}$ sont les deux
  projections de $\symup{R}$ dans $\symup{X}$;
  \item[(iii)] le morphisme $\symup{X/(p_1,p_2)\to Y}$ induit par $\symup{f}$ est un
  isomorphisme.
\end{enumerate}
Alors $\symup{f}$ est bien un épimorphisme, et même un épimorphisme strict (cf. [1],
A, \S~2. 3), la réciproque étant vraie si le produit fibré $\symup{X\times_Y X}$
existe. On dit aussi que l'objet quotient de $\symup{X}$ défini par l'épimorphisme $\symup{f}$
est un \uline{quotient effectif} de $\symup{X}$.

Les définitions précédentes impliquent la \fquote{\uline{correspondance galoisienne}}
suivante :

\src{102}
\result{PROPOSITION}{1.1} Il y a correspondance biunivoque, respectant les
ordres naturels, entre l'ensemble des relations d'équivalence effectives $\symup{R}$
dans $\symup{X}$, et l'ensemble des quotients effectifs $\symup{Y}$ de $\symup{X}$, à $\symup{R}$
correspondant le quotient effectif $\symup{X/R}$, et à $\symup{Y}$ correspondant la relation
d'équivalence effective définie par la projection canonique
$\symup{X\to Y}$, (qui est définie par le produit fibré $\symup{X\times_Y X}$
muni de ses deux projections).

Dans les très bonnes catégories (ensembles, faisceaux d'ensembles, etc.) tout
quotient est effectif, et toute relation d'équivalence est effective. Il n'en
est plus de même dans les catégories telles que la catégorie des préschémas
au-dessus d'un préschéma $\symup{S}$ donné, même lorsque $\symup{S}$ est le spectre d'un corps,
et même en se bornant aux schémas finis sur $\symup{S}$. Les questions d'effectivité,
et même (dans le cas de préschémas non finis sur $\symup{S}$) les questions d'existence
de quotients, s'avèrent le plus souvent délicates.

\subsection*{\textnormal{2. \uline{Exemple} : \uline{préschémas finis sur} $\symup{S}$.}}
\addcontentsline{toc}{subsection}{2. Exemple: préschémas finis sur S}

Soit $\symup{\text{\uwave{C}}}$ la catégorie des préschémas finis sur $\symup{S}$, supposé
localement noethérien. Elle est équivalente à la catégorie opposée de la
catégorie des faisceaux cohérents d'algèbres commutatives sur $\symup{S}$, ou encore,
si $\symup{S}$ est affine d'anneau $\symup{A}$, à la catégorie opposée de la catégorie des
$\symup{A}$-algèbres finies sur $\symup{A}$ (i. e. qui sont des modules de type fini sur $\symup{A}$).
On en conclut tout de suite que dans $\symup{\text{\uwave{C}}}$ les limites projectives
finies et les limites inductives finies existent. C'est bien connu (sans
aucune hypothèse de finitude\ldots) pour les premières. Ainsi le produit fibré
de préschémas $\symup{X,Y}$ sur $\symup{S}$ correspond au produit tensoriel $\symup{B\otimes_A C}$
des algèbres correspondantes, le noyau de deux morphismes
$\symup{X\rightrightarrows Y}$, défini par deux homomorphismes de $\symup{A}$-algèbres
$\symup{u,v:C\rightrightarrows B}$, correspond au quotient de $\symup{B}$ par l'idéal engendré
par les $\symup{u(c)-v(c)}$, etc. Pour les limites inductives finies, il suffit de
considérer d'une part les sommes finies, qui correspondent au produit ordinaire
de $\symup{A}$-algèbres, et d'autre part les conoyaux de couples de morphismes
$\symup{X\rightrightarrows Y}$, qui correspondent en effet, (comme on constate
aussitôt) au sous-anneau de $\symup{C}$ ensemble des éléments où les homomorphismes
$\symup{u,v:C\rightrightarrows B}$ coïncident (ce dernier est fini sur $\symup{A}$ grâce à
l'hypothèse noethérienne). Notons d'ailleurs qu'on peut montrer, utilisant
l'hypothèse noethérienne, que les limites inductives finies, et en particulier
les quotients, ainsi construits dans la catégorie $\symup{\text{\uwave{C}}}$ des préschémas
finis sur $\symup{S}$, sont en réalité des quotients dans la catégorie de \uline{tous}
les préschémas.

\src{103}
Comme nous avons dit dans [1], \uline{il y a dans} $\symup{\text{\uwave{C}}}$ \uline{des épimorphismes non effectifs} (ou stricts, cela revient au même puisque les produits fibrés
existent). \uline{J'ignore si les relations d'équivalence sont toujours
effectives}, lorsqu'on ne fait pas d'hypothèse de platitude. Je n'ai obtenu,
dans cette direction, que des résultats très partiels, positifs,
indispensables pour la démonstration du théorème fondamental en théorie
formelle des modules (cf. [2], B, th. 1). Signalons qu'il est facile dans le
problème posé de se ramener au cas où $\symup{S}$ est le spectre d'un anneau $\symup{A}$
artinien local, de corps résiduel algébriquement clos. Mais même si $\symup{A}$ est un
corps, la réponse n'est pas connue.

On peut aussi considérer le cas d'un préschéma $\symup{X}$ sur $\symup{S}$ qui n'est plus
supposé fini sur $\symup{S}$, mais en considérant une relation d'équivalence $\symup{R}$ dans
$\symup{X}$ telle que $\symup{p_1:R\to X}$ soit un morphisme fini. On dit alors
que $\symup{R}$ est une relation d'équivalence finie. Supposant pour simplifier $\symup{S}$ et
$\symup{X}$ affines (donc $\symup{R}$ affine, de sorte que la situation est ramenée à une
situation de pure algèbre commutative), \uline{on ignore même dans ce cas s'il
existe un quotient} $\symup{X/R=Y}$ \uline{et si le morphisme canonique}
$\symup{X\to Y}$ \uline{est fini}. (Le cas le plus simple est celui où on
suppose que $\symup{S}$ est le spectre d'un corps $\symup{k}$, et où $\symup{X}$ est le spectre de
$\symup{k[t]}$, i. e. la droite affine.) Bien entendu, si les deux problèmes précédents
se résolvent par l'affirmative, on peut conclure dans la situation présente que
$\symup{R}$ est effective. Notons que le problème de \uline{l'existence} d'un quotient
$\symup{Y}$ et de la finitude de $\symup{f:X\to Y}$ se pose dans exactement les
mêmes termes si, au lieu d'un graphe d'équivalence dans $\symup{X}$, on a seulement un
prégraphe d'équivalence dans $\symup{X}$, au sens du paragraphe 4.

La question du passage au quotient par une relation d'équivalence finie plus
ou moins arbitraire se pose dans la construction des préschémas par
\fquote{recollement} de préschémas donnés $\symup{X_i}$ suivant certains sous-préschémas
fermés; la loi du recollement s'exprime précisément par une relation
d'équivalence finie dans le préschéma somme $\symup{X}$ des $\symup{X_i}$. Il faut s'attendre
aussi que la solution des problèmes posés ici et de diverses variantes, sera
une condition préliminaire à la mise au point d'une technique générale de
constructions non projectives, dans la direction inaugurée dans [2].

Le seul fait général positif connu du rédacteur est le suivant :

\result{PROPOSITION}{2.1} Soient $\symup{S}$ un préschéma localement noethérien, $\symup{s}$
un point de $\symup{S}$, et $\symup{\Omega}$ une extension algébriquement close de $\symup{k(s)}$.
Considérons le
\src{104}\fquote{foncteur-fibre} correspondant $\symup{F}$, associant à tout $\symup{S}$-schéma
$\symup{X}$ fini sur $\symup{S}$, l'ensemble des points de $\symup{X/S}$ à valeurs dans $\symup{\Omega}$. Ce
foncteur (trivialement exact à gauche) est \uline{exact à droite}, i. e. commute
aux limites inductives finies, et en particulier aux conoyaux de couples de
morphismes.

En utilisant ce résultat pour tous les \fquote{points géométriques} de $\symup{S}$, on en
déduit que la catégorie $\symup{C'}$ \fquote{quotient} de $\symup{C}$, obtenue en
raisonnant \fquote{modulo morphismes surjectifs radiciels} (i. e. obtenue en
adjoignant formellement des inverses pour les morphismes en question), est une
catégorie \fquote{géométrique}, i. e. satisfait les mêmes propriétés \fquote{de nature
finie} que la catégorie des ensembles. En particulier, toute relation
d'équivalence y est effective. Cela implique que si $\symup{R}$ est une relation
d'équivalence dans un $\symup{X}$ fini sur $\symup{S}$, alors le morphisme canonique $\symup{R\to X\times_Y X}$ (où $\symup{Y=X/R}$) est \uline{radiciel et surjectif} (en fait une immersion fermée surjective,
puisque c'est un monomorphisme).

\subsection*{\textnormal{3. \uline{Cas d'un groupe d'opérateurs}.}}
\addcontentsline{toc}{subsection}{3. Cas d'un groupe d'opérateurs}

On suppose de nouveau que $\symup{\text{\uwave{C}}}$ est une catégorie quelconque. Soient
$\symup{G,X}$ des objets de $\symup{\text{\uwave{C}}}$, et supposons que $\symup{G}$ est un
$\symup{\text{\uwave{C}}}$-groupe d'opérateurs sur l'objet $\symup{X}$ de $\symup{\text{\uwave{C}}}$. Cela
signifie (cf. [2], A, \S~1) que pour tout objet $\symup{T}$ de $\symup{\text{\uwave{C}}}$, on s'est
donné une structure de groupe sur $\symup{G(T)}$ et une structure d'ensemble à groupe
d'opérateurs $\symup{G(T)}$ sur $\symup{X(T)}$, de telle façon que pour $\symup{T}$ variable, les
structures en question \fquote{varient fonctoriellement} avec $\symup{T}$. Si dans
$\symup{\text{\uwave{C}}}$ les produits $\symup{G\times G}$ et $\symup{G\times X}$ existent, une telle
structure peut encore se définir comme un couple de morphismes
\[
  \symup{G\times G\to G,
  \qquad \pi:G\times X\to X}
\]
soumis à la condition que pour tout objet $\symup{T}$ de $\symup{\text{\uwave{C}}}$, les lois de
composition correspondantes pour les ensembles $\symup{G(T)}$ et $\symup{X(T)}$ fassent de
$\symup{G(T)}$ un groupe opérant sur $\symup{X(T)}$. La traduction de cet axiome par la
commutativité de certains diagrammes dans $\symup{\text{\uwave{C}}}$ est facile, mais
fastidieuse, et en fait, parfaitement inutile dans tous les cas à ma
connaissance.

Supposons que $\symup{G\times X}$ existe, et considérons les deux morphismes
\[
  \symup{p_1,p_2:G\times X\rightrightarrows X}
\]
avec
\src{105}
\[
  \symup{p_1=\operatorname{pr}_1,\qquad p_2=\pi.}
\]
On constate aussitôt que le couple $\symup{(p_1,p_2)}$ est un couple d'équivalence si,
et seulement si, pour tout objet $\symup{T}$ de $\symup{\text{\uwave{C}}}$, l'application
\[
  \symup{G(T)\times X(T)\sim(G\times X)(T)
  \to X(T)\times X(T),}
\]
définie par ce couple, est injective, i. e. si le groupe $\symup{G(T)}$ opère
\uline{librement} sur l'ensemble $\symup{X(T)}$, i. e. $\symup{g\in G(T)}$, $\symup{x\in X(T)}$,
$\symup{g\cdot x=x}$ implique $\symup{g=}$ élément unité du groupe $\symup{G(T)}$. On dit
alors que $\symup{G}$ \uline{opère librement} sur $\symup{X}$, (ou que $\symup{X}$ est un
$\symup{\text{\uwave{C}}}$\uline{-espace principal sous} $\symup{G}$). La relation d'équivalence
associée au couple $\symup{(p_1,p_2)}$ est alors appelée la \uline{relation
d'équivalence définie par le groupe} $\symup{G}$ opérant librement sur $\symup{X}$. Lorsque
$\symup{X\times X}$ existe également et qu'on considère le morphisme
\[
  \symup{p:G\times X\to X\times X}
\]
défini par le couple $\symup{(p_1,p_2)}$, la condition que $\symup{G}$ opère librement signifie
que $\symup{p}$ est un \uline{monomorphisme}.

Bien entendu, même lorsque $\symup{G}$ n'opère pas librement sur $\symup{X}$, on désire avoir
des critères d'existence d'un quotient de $\symup{X}$ par $\symup{G}$, i. e. du conoyau du
couple $\symup{(p_1,p_2)}$ précédent.

Le conoyau en question sera souvent noté $\symup{X/G}$, ou de préférence $\symup{G\backslash
X}$ si $\symup{G}$ opère à gauche (la notation précédente étant réservée au cas où $\symup{G}$
opère à droite). On notera que même lorsque \fquote{l'image} de $\symup{G\times X}$ par $\symup{p}$
existe (cette image étant définie par exemple comme le plus petit sous-objet de
$\symup{X\times X}$ par lequel on peut factoriser $\symup{p}$), soit $\symup{R}$, cette dernière n'est
pas le plus souvent une relation d'équivalence dans $\symup{X}$. Si on essaye alors de
passer directement au quotient par $\symup{R}$ (ou plus précisément, par le couple des
morphismes de $\symup{R}$ dans $\symup{X}$ induits par les deux projections
$\symup{\operatorname{pr}_i}$) on perd sans espoir les caractères particuliers du
couple de départ $\symup{(p_1,p_2)}$. Il importe donc de trouver une généralisation de
la notion de relation d'équivalence, s'appliquant directement au couple défini
par un $\symup{\text{\uwave{C}}}$-groupe d'opérateurs.

\subsection*{\textnormal{4. \uline{Pré-relations d'équivalence}.}}
\addcontentsline{toc}{subsection}{4. Pré-relations d'équivalence}

Rappelons qu'on appelle \uline{groupoïde} une catégorie où tous les morphismes
sont des isomorphismes. D'autre part, une catégorie doit être définie comme
formée de deux ensembles de base $\symup{(X,R)}$, l'ensemble des \uline{objets} et
l'ensemble des \uline{flèches}, munis des structures suivantes :
\src{106}
\begin{enumerate}
\item[(i)] Un couple d'applications
\[
  \symup{p_1,p_2:R\rightrightarrows X}
\]
appelées \uline{application-source} et \uline{application-but},

\item[(ii)] Une application
\[
  \symup{\pi:(R,p_2)\times_X(R,p_1)\to R}
\]
appelée \uline{application-composition}.
\end{enumerate}
Ces données doivent satisfaire à des axiomes bien connus, que nous ne
répéterons pas ici, et qui pourraient s'expliciter par la commutativité de
certains diagrammes, et l'existence d'une application $\symup{D:X\to R}$
(nécessairement unique) rend commutatifs deux autres diagrammes. $\symup{D}$
correspond au passage d'un objet à la flèche identique correspondante, et
satisfait à
\[
  \symup{p_1D=p_2D=\operatorname{id}_X.}
\]
Dire que la catégorie est un groupoïde signifie alors qu'il existe une
application
\[
  \symup{s:R\to R}
\]
(nécessairement unique), appelée la \uline{symétrie} de $\symup{R}$, transformant
chaque flèche en une flèche inverse, ce qui pourrait s'expliciter par la
commutativité de quatre autres diagrammes, formés au moyen de $\symup{s}$, $\symup{\Delta}$
et des données précédentes, et dont les deux premiers s'écrivent :
\[
  \symup{p_1s=p_2,\qquad p_2s=p_1.}
\]

Ces points étant rappelés, les définitions générales dans ([2], A,
\S~1) montrent en particulier ce qu'il faut entendre par une structure de
$\symup{\text{\uwave{C}}}$-\uline{catégorie}, resp. de $\symup{\text{\uwave{C}}}$-\uline{groupoïde}, sur un
couple d'objets $\symup{X,R}$ d'une catégorie quelconque $\symup{\text{\uwave{C}}}$ : c'est, par
définition, la donnée, pour tout objet $\symup{T}$ de $\symup{\text{\uwave{C}}}$, d'une structure
de catégorie resp. de groupoïde au sens ensembliste, dont l'ensemble des
objets soit $\symup{X(T)}$, et l'ensemble des flèches $\symup{R(T)}$, ces structures
\fquote{variant fonctoriellement} avec $\symup{T}$ variable.
Cela implique donc la définition de deux morphismes
\[
  \symup{p_1,p_2:R\rightrightarrows X}
\]
appelés \uline{morphisme-source} et \uline{morphisme-but}, et lorsque le produit
fibré voulu existe d'un morphisme
\src{107}
\[
  \symup{\pi:(R,p_2)\times_X(R,p_1)\to R}
\]
appelé morphisme composition, ces trois données suffisant alors à
déterminer la structure de catégorie sur $\symup{X,R}$, l'axiome à mettre sur ces
données étant le suivant : pour tout $\symup{T}$, les trois données correspondantes
pour $\symup{X(T),R(T)}$ définissent sur ce couple d'ensembles une structure de
catégorie, (resp. de groupoïde). Le cas échéant, cela peut s'exprimer par la
commutativité de certains diagrammes, impliquant un morphisme bien déterminé
\[
  \symup{D:X\to R}
\]
et dans le cas des groupoïdes, un morphisme bien déterminé
\[
  \symup{s:R\to R,}
\]
diagrammes qui s'explicitent comme dans le cas \fquote{ensembliste}. Cette
interprétation fastidieuse des axiomes est heureusement inutile pour la
pratique, le seul intérêt théorique de la possibilité d'exprimer les données
et axiomes à partir de morphismes et égalités de morphismes entre certains
produits fibrés étant le suivant : si on a un foncteur
$\symup{F:\text{\uwave{C}}\to\text{\uwave{C}}'}$ exact à gauche (i. e. commutant aux produits
finis et aux produits fibrés), il transforme une $\symup{\text{\uwave{C}}}$-catégorie
(resp. un $\symup{\text{\uwave{C}}}$-groupoïde) en une $\symup{\text{\uwave{C}}'}$-catégorie (resp. en un
$\symup{\text{\uwave{C}}'}$-groupoïde) (sous réserve de l'existence des produits finis et
produits fibrés dans $\symup{\text{\uwave{C}}}$).

Il est important en pratique de savoir interpréter les morphismes
$\symup{p_1,p_2,\pi,D,s}$ comme des \uline{opérations simpliciales} dans un objet
semi-simplicial ou simplicial convenable de $\symup{\text{\uwave{C}}}$, du moins lorsque
dans $\symup{\text{\uwave{C}}}$ les produits fibrés existent. Pour fixer la terminologie,
introduisons la \uline{catégorie} $\symup{\text{\uwave{S}}}$ \uline{des simplexes-types} comme la
catégorie dont les objets sont les ensembles finis de la forme
\[
  \symup{\Delta_n=[0,n]\qquad\text{où }n\geqslant 0}
\]
(intervalle des entiers de $\symup{0}$ à $\symup{n}$), et dont les morphismes ou flèches sont
les \uline{applications quelconques} entre ces ensembles finis. On notera que
la catégorie $\symup{\text{\uwave{S}}}$ est équivalente à la catégorie des ensembles finis
\uline{non vides}, où on prend comme morphismes les applications entre ensembles
finis. Dans $\symup{\text{\uwave{S}}}$, la somme d'une famille finie \uline{non vide} d'objets
existe évidemment, ainsi que la somme amalgamée de deux objets sous un
troisième (opération duale du produit fibré). On désigne par $\symup{\text{\uwave{\phantom{S}}}'}$ la
sous-catégorie de $\symup{\text{\uwave{S}}}$ ayant mêmes objets, mais où les morphismes sont
les \uline{applications croissantes} entre les $\symup{\Delta_n}$. Cette catégorie est
équivalente à la catégorie des ensembles finis totalement ordonnés non
vides. Dans cette catégorie la
\src{108}
somme de deux objets n'existe jamais, et la somme amalgamée de deux objets
$\symup{A,B}$ sous un troisième $\symup{C}$ n'existe pas en général (prendre par exemple
$\symup{C=\Delta_0}$, $\symup{A=B=\Delta_1}$, les deux applications structurales
$\symup{u:C\to A}$ et $\symup{v:C\to B}$ égales). Elle existe cependant dans certains cas,
par exemple
\[
  \symup{C=\Delta_0,\qquad A=\Delta_m,\qquad B=\Delta_n,\qquad
  u(0)=m,\qquad v(0)=0}
\]
et dans ce cas, on a
\[
  \symup{\Delta_m\mathbin{\mathrm{n}}_{\Delta_0}\Delta_n=\Delta_{m+n}.}
\]

Un \uline{objet simplicial} (resp. un objet \uline{semi-simplicial}) dans une
catégorie $\symup{\text{\uwave{C}}}$ est par définition un foncteur contravariant $\symup{K}$ de
$\symup{\text{\uwave{S}}}$ (resp. de $\symup{\text{\uwave{S}}'}$) dans $\symup{\text{\uwave{C}}}$. Un objet simplicial
définit donc un objet semi-simplicial par restriction, mais la première
notion diffère de la seconde essentiellement par la présence d'\uline{opérations
de symétrie} dans les $\symup{K_n=K(\Delta_n)}$, qui correspondent aux transformés
par le foncteur $\symup{K}$ des éléments du groupe symétrique à $\symup{n+1}$ lettres
(considéré comme le groupe des automorphismes de $\symup{\Delta_n}$ dans
$\symup{\text{\uwave{S}}}$).

Ceci dit, pour tout $\symup{n}$, soit $\symup{\Delta'_n}$ (resp. $\symup{\Delta''_n}$) la catégorie
finie, dont l'ensemble des objets est $\symup{\Delta_n}$, définie par la relation
d'ordre chaotique (resp. d'ordre total naturel) sur $\symup{\Delta_n}$ (l'ensemble
des flèches étant le graphe de la dite relation d'ordre). Il est évident que
$\symup{\Delta'_n}$ (resp. $\symup{\Delta''_n}$) dépend fonctoriellement de l'objet $\symup{\Delta_n}$
de $\symup{\text{\uwave{S}}}$ (resp. $\symup{\text{\uwave{S}}'}$). Si alors $\symup{Z}$ est une catégorie,
$\symup{\operatorname{Hom}(\Delta'_n,Z)}$ (resp.
$\symup{\operatorname{Hom}(\Delta''_n,Z)}$) est, pour $\symup{\Delta_n}$ variable, un foncteur
de la catégorie $\symup{\text{\uwave{S}}}$, (resp. $\symup{\text{\uwave{S}}'}$), dans la catégorie des
ensembles, i. e. un \uline{ensemble simplicial} (resp. \uline{semi-simplicial}), dit
\uline{associé à la catégorie} $\symup{Z}$ : soient $\symup{Z'}$ et $\symup{Z''}$. On a d'ailleurs un
homomorphisme naturel évident de l'ensemble semi-simplicial associé à $\symup{Z'}$
dans $\symup{Z''}$, et ce dernier est un isomorphisme si et seulement si $\symup{Z}$ est un
groupoïde. Ceci posé :

\result{PROPOSITION}{4.1} Le foncteur $\symup{Z\rightsquigarrow Z''}$ de la catégorie des
catégories dans la catégorie des ensembles semi-simpliciaux est pleinement
fidèle, et définit une équivalence de la catégorie des \uline{catégories} avec
la catégorie des ensembles semi-simpliciaux, i. e. des foncteurs contravariants
$\symup{K}$ de $\symup{\text{\uwave{S}}'}$ dans $\symup{(\mathrm{Ens})}$, \uline{qui transforment les sommes amalgamées} $\symup{A\mathbin{\phantom{\mathrm{n}}}_C B}$ \uline{du type précisé plus haut, en produits fibrés d'ensembles}. De même, le foncteur $\symup{Z\rightsquigarrow Z'}$ de la catégorie des
groupoïdes dans la catégorie des ensembles simpliciaux est pleinement fidèle,
et définit une équivalence de la catégorie des \uline{groupoïdes} avec la
catégorie des
\src{109}
ensembles simpliciaux, i. e. des foncteurs contravariants $\symup{K}$ de $\symup{\text{\uwave{S}}}$
dans $\symup{(\mathrm{Ens})}$, \uline{qui transforment des sommes amalgamées en produits fibrés}.

On peut donc considérer les catégories comme des ensembles semi-simpliciaux
particuliers, et les groupoïdes comme des ensembles simpliciaux particuliers,
à condition bien entendu de raisonner \fquote{à un isomorphisme près} comme il est
de rigueur quand on interprète certaines structures en termes d'autres. Le
procédé habituel de réduction au cas ensembliste implique alors :

\result{COROLLAIRE}{4.2} L'énoncé précédent reste valable quand on remplace
les catégories, groupoïdes, ensembles simpliciaux, par des
$\symup{\text{\uwave{C}}}$-catégories, $\symup{\text{\uwave{C}}}$-groupoïdes, $\symup{\text{\uwave{C}}}$-objets
simpliciaux, pourvu que dans $\symup{\text{\uwave{C}}}$ les produits fibrés existent.

L'objet semi-simplicial $\symup{K}$ dans $\symup{\text{\uwave{C}}}$ associé à une catégorie
$\symup{(X,R,\ldots)}$ dans $\symup{\text{\uwave{C}}}$ peut s'expliciter, en interprétant la
composante $\symup{K_n=K(\Delta_n)}$ de $\symup{K}$ comme étant la puissance fibrée
$\symup{(n+1)}$-ième de $\symup{(R,p_1)}$ sur $\symup{X}$, ou mieux,
par la formule de récurrence,
\[
  \symup{K_0=R,\qquad
  K_n=(K_{n-1},p_n^{(n-1)})\times_X(R,p_1)}
\]
où les $\symup{p_i^{(n-1)}}$ ($\symup{0\leqslant i\leqslant n-1}$) sont les projections naturelles de
$\symup{K_{n-1}}$ dans $\symup{X}$ (qui se définissent également par récurrence). De cette
façon, $\symup{p_1,p_2,\pi,D,s}$ s'interprètent comme les opérations simpliciales qui
correspondent respectivement aux morphismes dans $\symup{\text{\uwave{S}}}$ : face $\symup{0}$ de
$\symup{\Delta_1}$, face $\symup{1}$ de $\symup{\Delta_1}$, face $\symup{(0,2)}$ de $\symup{\Delta_2}$,
dégénérescence $\symup{\Delta_1\to\Delta_0}$, symétrie de $\symup{\Delta_1}$. Toutes les
autres opérations semi-simpliciales (resp. simpliciales) se déduisent
formellement des quatre (resp. cinq) précédentes par composition et produits
fibrés.

On appelle maintenant \uline{pré-relation d'équivalence} dans un objet $\symup{X}$
d'une catégorie, la donnée d'un groupoïde dont l'objet des objets (si on peut
dire) est $\symup{X}$. Une telle donnée comporte donc entre autres un objet $\symup{R}$ et
deux morphismes
\[
  \symup{p_1,p_2:R\rightrightarrows X.}
\]
Mais on notera que ces seules données ne déterminent pas la structure
envisagée, contrairement à ce qui a lieu pour les couples d'équivalence. Dans
cet exposé, nous nous intéressons à cette notion dans le but d'obtenir des
critères de possibilité de passage au quotient, i. e. de formation du conoyau
du couple $\symup{(p_1,p_2)}$. L'énoncé du problème ne fait donc pas appel aux données
supplémentaires inhérentes à un groupoïde. Dans la démonstration des résultats
qui vont suivre, on se
\src{110}
sert cependant de ces données supplémentaires, et en particulier des
opérations simpliciales (y inclus des opérations de symétrie) jusqu'en
dimension $\symup{3}$ (faisant intervenir la puissance fibrée quadruple de $\symup{R}$ sur
$\symup{X}$).

Une relation d'équivalence dans un objet $\symup{X}$ de $\symup{\text{\uwave{C}}}$ définit une
pré-relation d'équivalence : il suffit de le voir dans le cas ensembliste, et
on associe alors à une relation d'équivalence, dans un ensemble $\symup{X}$, le
groupoïde dont l'ensemble d'objets est $\symup{X}$, et l'ensemble des flèches est
l'ensemble graphe de la relation d'équivalence.

Un $\symup{\text{\uwave{C}}}$-monoïde $\symup{G}$ opérant sur un objet $\symup{X}$ de $\symup{\text{\uwave{C}}}$ définit
une $\symup{\text{\uwave{C}}}$-catégorie dont les objets de base sont $\symup{R=G\times X}$ et $\symup{X}$
(sous réserve que $\symup{G\times X}$ existe), et qui est un $\symup{\text{\uwave{C}}}$-groupoïde
si et seulement si $\symup{G}$ est un groupe. Il suffit encore d'en faire la
vérification dans le cas ensembliste. On définit alors la composition des
flèches $\symup{(g,a)}$ et $\symup{(g',g\cdot a)}$ comme étant
\[
  \symup{(g',g\cdot a)(g,a)=(g'g,a)}
\]
i. e. si $\symup{a,b\in X}$ ; alors $\symup{\operatorname{Hom}(a,b)}$ est par définition le
transporteur de $\symup{a}$ dans $\symup{b}$, et les morphismes se composent grâce à la
composition des éléments de $\symup{G}$.

\result{REMARQUE}{} On peut éviter les difficultés logiques que
soulève un énoncé tel que (4.1) en y sous-entendant que tous les objets
envisagés se trouvent dans un \fquote{Univers} fixé, (qui lui-même est un
ensemble).

\fsection{\normalfont 5}{\normalfont\uline{Quotient par une relation d'équivalence finie et plate}.}

\result{THÉORÈME}{5.1} Soient $\symup{X=\operatorname{Spec}(B)}$ un schéma affine,
$\symup{\mathcal R}$ une pré-relation d'équivalence dans $\symup{X}$, dont la composante $\symup{R_1}$ est
affine, soit $\symup{R_1=\operatorname{Spec}(C)}$. On suppose que la première
projection $\symup{p_1:R_1\to X}$ est un morphisme fini et localement libre, en
d'autres termes que l'homomorphisme d'anneaux correspondant
$\symup{p_1':B\to C}$ fait de $\symup{C}$ un $\symup{B}$-module projectif de type fini. Soit $\symup{A}$ le
sous-anneau de $\symup{B}$ noyau du couple d'homomorphismes $\symup{p_1',p_2':B\rightrightarrows C}$ (= ensemble des éléments $\symup{b}$ tels que $\symup{p_1'(b)=p_2'(b)}$). Soit
$\symup{Y=\operatorname{Spec}(A)}$, et $\symup{f:X\to Y}$ le morphisme défini par l'immersion
de $\symup{A}$ dans $\symup{B}$. Sous ces conditions :
\begin{enumerate}
\item[(i)] $\symup{B}$ est entier sur $\symup{A}$, i. e. $\symup{f}$ est un morphisme entier.

\item[(ii)] Le morphisme $\symup{f}$ est surjectif, ses fibres sont les classes
d'équivalence
\src{111}
ensemblistes $\symup{p_2\allowbreak(p_1^{-1}(x))}$ dans $\symup{X}$ modulo $\symup{\mathcal R}$, et la topologie de $\symup{Y}$
est quotient de celle de $\symup{X}$.

\item[(iii)] $\symup{Y}$ est le quotient de $\symup{X}$ par $\symup{\mathcal R}$ dans la catégorie des
préschémas.

\item[(iv)] Si $\symup{\mathcal R}$ provient d'une \uline{relation} d'équivalence, alors le
morphisme $\symup{f:X\to Y}$ est fini localement libre (i. e. $\symup{B}$ est un $\symup{A}$-module
projectif de type fini), et la relation d'équivalence est effective, i. e.
$\symup{R_1\to X\times_Y X}$ est un isomorphisme.
\end{enumerate}

Ce théorème généralise le théorème bien connu relatif au cas d'un groupe
fini $\symup{G}$ opérant par automorphismes sur l'anneau $\symup{B}$, et à l'anneau $\symup{A}$ des
invariants, et la démonstration est analogue à la démonstration connue. On
peut préciser (iii) ainsi :

\result{COROLLAIRE}{5.2} Le morphisme canonique
$\symup{R_1\to X\times_Y X}$ est \uline{surjectif}.

Soit toujours $\symup{\mathcal R}$ une pré-relation d'équivalence \fquote{finie et localement libre} dans $\symup{X}$, mais $\symup{X}$ étant un préschéma quelconque. Supposons qu'on
puisse trouver un préschéma $\symup{Y}$ et un morphisme $\symup{f:X\to Y}$ tel que
$\symup{fp_1=fp_2}$, et tel en plus que la suite d'homomorphismes de faisceaux
d'anneaux sur $\symup{Y}$ :
\[
  \symup{\mathcal O_Y\to f_*(\mathcal O_X)
  \rightrightarrows g_*(\mathcal O_R)}
\]
soit exacte (où $\symup{g=fp_i}$). Il résulte alors du théorème qu'on a des
conclusions (i) à (iv) analogues à celles du théorème, en particulier, par
(iii), $\symup{Y}$ est le quotient de $\symup{X}$ par $\symup{\mathcal R}$, et par suite déterminé à un
isomorphisme unique près. On dira sous ces conditions que la pré-relation
d'équivalence $\symup{\mathcal R}$ dans $\symup{X}$ est \fquote{\uline{admissible}}. Avec cette définition :

\result{THÉORÈME}{5.3} Soient $\symup{X}$ un préschéma, $\symup{\mathcal R}$ une pré-relation
d'équivalence dans $\symup{X}$, telle que $\symup{p_1:R_1\to X}$ soit un morphisme fini et
localement libre. Pour que $\symup{\mathcal R}$ soit admissible, il faut et il suffit que toute
classe d'équivalence ensembliste $\symup{p_2(p_1^{-1}(x))}$ dans $\symup{X}$ modulo $\symup{\mathcal R}$ soit
contenue dans un ouvert affine (condition toujours vérifiée si toute partie
finie de $\symup{X}$ est contenue dans un ouvert affine, par exemple si $\symup{X}$ est
quasi-projectif sur un schéma affine).

On montre en effet sans difficulté que toute classe d'équivalence mod $\symup{\mathcal R}$
dans $\symup{X}$ est alors contenue dans un ouvert affine \uline{stable} par $\symup{\mathcal R}$, et on
construit le quotient $\symup{Y}$ par recollement des morceaux obtenus à l'aide du
théorème 5.1.
\src{112}

\result{COROLLAIRE}{5.4} Supposons cette condition réalisée et que de plus
$\symup{\mathcal R}$ provient d'une relation d'équivalence. Alors cette dernière est effective,
i. e. $\symup{R_1\to X\times_Y X}$ est un isomorphisme, et $\symup{f:X\to Y}$ est un morphisme
fini localement libre.

On en conclut aussitôt, par \fquote{descente} :

\result{COROLLAIRE}{5.5} Sous les conditions de (5.4), pour que $\symup{X}$ soit
partout de rang $\symup{n}$ au-dessus de $\symup{Y}$, il faut et il suffit que $\symup{(R_1,p_1)}$ soit
partout de rang $\symup{n}$ au-dessus de $\symup{X}$. Si $\symup{X,R_1}$ sont des $\symup{Z}$-préschémas et
$\symup{p_1,p_2}$ des $\symup{Z}$-morphismes, donc $\symup{Y}$ un $\symup{Z}$-préschéma, alors $\symup{X}$ est plat
sur $\symup{Z}$ si et seulement si $\symup{Y}$ l'est.

En résumé :

\result{SCHOLIE}{} La donnée d'un morphisme fini, localement libre et
surjectif $\symup{f:X\to Y}$ de préschémas, est équivalente à la donnée d'un
préschéma $\symup{X}$ muni d'une relation d'équivalence $\symup{R}$ telle que $\symup{p_1:R\to X}$
soit fini et localement libre, et que toute classe $\symup{p_2(p_1^{-1}(x))}$ soit
contenue dans un ouvert affine.

\medskip
\noindent\textnormal{REMARQUES 5.6.}

\smallskip
\noindent $\symup{1^\circ}$ On n'a pas eu à faire d'hypothèse noethérienne.

\smallskip
\noindent $\symup{2^\circ}$ Cette notion de passage au quotient comprend comme cas
particulier la \fquote{descente inséparable} de CARTIER, qui correspond à la
détermination des morphismes finis localement libres $\symup{f:X\to Y}$ tels que
$\symup{f_*(\mathcal O_X)}$ admette une $\symup{p}$-base par rapport à $\symup{\mathcal O_Y}$ ($\symup{X}$
étant un préschéma donné dont le faisceau $\symup{\mathcal O_X}$ est annulé par le
nombre premier $\symup{p>0}$). On notera que le résultat s'exprime aisément sans
hypothèse de régularité sur les anneaux locaux et sans supposer que $\symup{X}$ est un
schéma algébrique sur un corps. La théorie de JACOBSON-BOURBAKI est obtenue
en prenant pour $\symup{X}$ le spectre d'un corps de caractéristique $\symup{p}$.

\smallskip
\noindent $\symup{3^\circ}$ GABRIEL avait auparavant obtenu un cas particulier du
théorème (5.3) dans la théorie du passage au quotient pour les groupes
commutatifs finis au-dessus d'un corps $\symup{k}$. [Comparer (7.4)].

\fsection{\normalfont 6}{\normalfont\uline{Quotient par une relation d'équivalence propre et plate}.}

\result{THÉORÈME}{6.1} Soient $\symup{S}$ un préschéma \uline{localement noethérien}, $\symup{X}$
un $\symup{S}$-schéma \uline{quasi-projectif}, $\symup{\mathcal R}$ une pré-relation d'équivalence dans $\symup{X}$,
telle que :
\src{113}
(a) $\symup{p_1:R_1\to X}$ est propre et plat ; (b) $\symup{R_1\to X\times_S X}$ est un morphisme fini, ou ce qui revient au
même en vertu de (a), à fibres finies (condition automatiquement vérifiée si
$\symup{\mathcal R}$ provient d'une relation d'équivalence). Sous ces conditions :
\begin{enumerate}
\item[(i)] $\symup{X/\mathcal R=Y}$ existe, et (si $\symup{S}$ est noethérien) est quasi-projectif sur $\symup{S}$.

\item[(ii)] Le morphisme canonique $\symup{f:X\to Y}$ est surjectif, propre, ouvert,
ses fibres sont les classes d'équivalence $\symup{p_2(p_1^{-1}(x))}$ dans $\symup{X}$ mod $\symup{\mathcal R}$,
donc $\symup{Y}$ s'identifie à l'espace topologique quotient de $\symup{X}$ par la relation
d'équivalence ensembliste définie par $\symup{\mathcal R}$. Enfin,
$\symup{R_1\to X\times_Y X}$ est surjectif.

\item[(iii)] Si $\symup{\mathcal R}$ provient d'une relation d'équivalence, cette dernière est
effective, i. e. $\symup{R_1\to X\times_Y X}$ est un isomorphisme, et de plus
$\symup{f:X\to Y}$ est plat (donc fidèlement plat).
\end{enumerate}

Pour la démonstration, on se ramène à (5.1) grâce à des quasi-sections
convenables de $\symup{X}$ pour $\symup{\mathcal R}$, la démonstration étant analogue à la construction
des groupes algébriques quotients dans le Séminaire Chevalley.

En résumé :

\result{SCHOLIE}{} Soit $\symup{X}$ quasi-projectif sur $\symup{S}$ localement noethérien. La
donnée d'un morphisme propre, fidèlement plat et surjectif $\symup{f:X\to Y}$ de $\symup{X}$
dans un $\symup{S}$-préschéma $\symup{Y}$, est équivalente à la donnée d'une relation
d'équivalence $\symup{R}$ sur $\symup{X}$, telle que $\symup{p_1:R\to X}$ soit propre et plat.

La même méthode donne le résultat suivant :

\result{THÉORÈME}{6.2} Soient $\symup{S}$ un préschéma noethérien, $\symup{X}$ un
préschéma de type fini sur $\symup{S}$, $\symup{\mathcal R}$ une pré-relation d'équivalence dans le
$\symup{S}$-préschéma $\symup{X}$, on suppose que : (a) $\symup{p_1:R_1\to X}$ est plat et de type fini ; (b) Le morphisme $\symup{R_1\to X\times_S X}$ est quasi-fini, (i. e. à fibres
finies). Alors il existe un ouvert $\symup{U}$ \uline{dense} dans $\symup{X}$ \uline{saturé} pour $\symup{\mathcal R}$, tel que :
\begin{enumerate}
\item[(i)] Si $\symup{\mathcal R_U}$ est la pré-relation d'équivalence induite par $\symup{\mathcal R}$ dans
$\symup{U}$, alors $\symup{U/\mathcal R_U}$ existe, et est de type fini sur $\symup{S}$.

\item[(ii)] Le morphisme canonique $\symup{U\to U/\mathcal R_U}$ est surjectif et ouvert, ses
fibres sont les classes d'équivalence ensembliste pour $\symup{\mathcal R_U}$, (donc $\symup{U/\mathcal R_U}$ est
un espace topologique quotient de $\symup{U}$ par la relation d'équivalence ensembliste
définie par $\symup{\mathcal R_U}$), enfin le morphisme
$\symup{(\mathcal R_U)_1\to U\times_{U/\mathcal R_U}U}$ est surjectif.
\src{114}

\item[(iii)] Si $\symup{\mathcal R}$ provient d'une relation d'équivalence, on peut supposer
$\symup{U\to U/\mathcal R_U}$ fidèlement plat et $\symup{\mathcal R_U}$ effectif.
\end{enumerate}

C'est là un résultat de nature essentiellement \fquote{birationnelle}.

\medskip
\noindent\textnormal{REMARQUES 6.3.}

\smallskip
\noindent $\symup{1^\circ}$ J'ignore si dans (6.1) et (6.2) l'hypothèse (b) est
inutile. Elle nous oblige pratiquement dans le passage au quotient par des
groupes de nous borner au cas où les stabilisateurs sont tous des groupes
finis.

\smallskip
\noindent $\symup{2^\circ}$ On peut se demander s'il n'y a pas des résultats analogues
à (6.1) et (6.2) sans hypothèse de platitude. Je n'ai aucun contre-exemple
dans cette direction. Par contre, même en gardant l'hypothèse de platitude, et
en se bornant à des relations d'équivalence telles que $\symup{p_1:R\to X}$ soit plat
et quasi-fini (mais non fini), $\symup{X}$ étant affine, il peut arriver que $\symup{R}$ ne
soit pas effective : prendre les relations d'équivalence induites sur des
ouverts affines recouvrant la variété de Nagata (où un groupe à deux éléments
opère de façon \fquote{non admissible}).

\fsection{\normalfont 7}{\normalfont\uline{Applications}.}

Comme nous l'avons dit dans l'introduction, l'application la plus importante
de (6.1) est la construction de schémas de Picard, ainsi que de solutions pour
divers autres problèmes de \fquote{modules}, sur lesquels nous reviendrons
ultérieurement.

On obtient une démonstration simple du résultat suivant de SHIMURA :

\result{PROPOSITION}{7.1} Soit $\symup{A}$ un schéma abélien défini sur un anneau de
valuation discrète $\symup{V}$ de corps des fractions $\symup{K}$. Alors tout schéma abélien
$\symup{B'}$ sur $\symup{K}$ isogène à un quotient de $\symup{A\otimes_V K}$ \fquote{se réduit bien pour $\symup{V}$} i. e. est isomorphe à un $\symup{B\otimes_V A}$, où $\symup{B}$ est un
schéma abélien sur $\symup{V}$ (essentiellement unique, rappelons-le).

On peut supposer que $\symup{B'}$ est le quotient de $\symup{A_K}$ par un sous-schéma en
groupes $\symup{C'}$. (N. B. : $\symup{C'}$ ne sera en général pas \fquote{réduit}, i. e. ses anneaux
locaux auront des éléments nilpotents). Considérons le sous-schéma fermé $\symup{C}$
de $\symup{A}$ \fquote{adhérence} de $\symup{C'}$, i. e. le plus petit sous-schéma fermé de $\symup{A}$ tel
que $\symup{C_K}$ majore $\symup{C'}$. On aura $\symup{C_K=C'}$, et du fait que $\symup{V}$ est un anneau de
valuation discrète, on tire facilement que $\symup{C}$ est un sous-schéma en groupes
de $\symup{A}$ sur $\symup{V}$. Comme $\symup{A}$ est propre sur $\symup{\operatorname{Spec}(V)=S}$, il en est
de même de $\symup{C}$, d'autre part $\symup{A}$ est même projectif sur $\symup{S}$. On peut alors
appliquer (6.1) pour construire $\symup{A/C=B}$,
\src{115}
c'est le $\symup{B}$ cherché.

Enfin, des raisonnements essentiellement connus permettent de tirer de (6.2)
le résultat suivant :

\result{THÉORÈME}{7.2} Soient $\symup{S}$ le spectre d'un anneau artinien, $\symup{F}$ et
$\symup{G}$ des schémas en groupes de type fini sur $\symup{S}$, $\symup{u:F\to G}$ un homomorphisme
de schémas en groupes sur $\symup{S}$. On suppose
\begin{enumerate}
\item[(i)] $\symup{F}$ est plat sur $\symup{S}$;
\item[(ii)] le noyau de $\symup{u}$ est fini.
\end{enumerate}
Sous ces conditions, le schéma quotient $\symup{G/F}$ existe, le morphisme canonique
$\symup{G\to G/F}$ est surjectif et ouvert, ses fibres sont les classes d'équivalence
ensemblistes définies par les opérations à droite de $\symup{F}$ sur $\symup{G}$. Enfin, si
$\symup{u}$ est un monomorphisme, alors le morphisme $\symup{G\to G/F}$ est plat et le
morphisme $\symup{G\times F\to G\times_{(G/F)}G}$ est un isomorphisme, en d'autres termes $\symup{G}$ est un espace principal homogène
sur $\symup{G/F}$, de groupe structural (opérant à droite) $\symup{F}$, ou plutôt
$\symup{F\times_S(G/F)}$ considéré comme schéma en groupes sur $\symup{G/F}$ (cf. [1], B,
\S~6).

\result{COROLLAIRE}{7.3} Sous ces conditions, pour que $\symup{G}$ soit plat sur $\symup{S}$,
il faut et il suffit que $\symup{G/F}$ le soit. Si cette condition est vérifiée, alors
le passage au quotient par $\symup{F}$ commute à toute extension de la base $\symup{S}$, et
si $\symup{F}$ est un sous-groupe invariant de $\symup{G}$, $\symup{G/F}$ peut être muni d'une
structure de \uline{groupe quotient} de $\symup{G}$ par $\symup{F}$.

La situation est particulièrement simple si $\symup{S}$ est le spectre d'un corps, car
alors tout $\symup{S}$-préschéma est automatiquement plat sur $\symup{S}$. On trouve :

\result{COROLLAIRE}{7.4} Soient $\symup{F,G}$ des schémas en groupes de type fini sur
un corps $\symup{k}$, et soit $\symup{u:F\to G}$ un homomorphisme de $\symup{k}$-groupes. Alors $\symup{u}$ se
factorise en $\symup{F\to F'\to G}$, où $\symup{F\to F'}$ est un homomorphisme de passage au
quotient par le sous-groupe fermé $\symup{\operatorname{Ker}u}$ de $\symup{F}$, et où
$\symup{F'\to G}$ est un homomorphisme de groupes qui est une immersion fermée. Le
quotient $\symup{G/F=G/F'}$ existe. Le formalisme habituel (genre théorèmes de Noether)
est valable parmi les groupes algébriques sur $\symup{k}$.

Ce résultat permet de traiter de manière uniforme le passage au quotient dans
les groupes algébriques au sens classique (i. e. qui sont irréductibles sur $\symup{k}$
et simples sur $\symup{k}$), et le passage au quotient par des sous-groupes
\fquote{infinitésimaux}
\src{116}
considéré par CARTIER. Il y a avantage à considérer les \fquote{hyperalgèbres}
introduites par CARTIER, à la suite des travaux de DIEUDONNÉ sur les groupes
formels, comme des groupes dans la catégorie des schémas formels au-dessus de
$\symup{k}$, et le cas échéant (s'ils correspondent à des hyperalgèbres de rang fini
sur $\symup{k}$) comme des groupes algébriques finis sur $\symup{k}$.

\fsection{\normalfont 8}{\normalfont\uline{Une conjecture}.}

Elle répond au besoin de savoir passer au quotient par le groupe projectif
opérant sur certains sous-schémas des \fquote{schémas de Hilbert} (ces derniers
remplaçant, en théorie des schémas, les variétés de Chow).

Soient $\symup{S}$ un préschéma, et $\symup{n}$ un entier. A tout préschéma $\symup{S'}$ sur $\symup{S}$,
faisons correspondre le groupe
$\symup{\operatorname{Gl}(n,\Gamma(S',\mathcal O_{S'}))}$ des matrices inversibles à
$\symup{n}$ lignes et $\symup{n}$ colonnes à valeurs dans l'anneau des sections de
$\symup{\mathcal O_{S'}}$. On obtient ainsi un foncteur contravariant en $\symup{S'}$, dont on
montre aisément qu'il est représentable, il correspond donc à un schéma en
groupes sur $\symup{S}$, d'ailleurs affine sur $\symup{S}$, noté $\symup{G\ell(n)_S}$. Sa
formation est compatible avec le changement de base, de sorte qu'en réalité
tout provient d'un schéma en groupes sur $\symup{\text{\uwave{Z}}}$, noté $\symup{G\ell(n)}$. Le groupe $\symup{\operatorname{Gl}(1)}$, appelé groupe
multiplicatif et souvent noté $\symup{G_m}$, correspond au foncteur
$\symup{S\rightsquigarrow\Gamma(S,\mathcal O_S)^*}$, groupe des \fquote{unités} sur $\symup{S}$. On a
un homomorphisme évident $\symup{\operatorname{Gl}(1)\to\operatorname{Gl}(n)}$, et on
construit facilement le groupe quotient, noté $\symup{\operatorname{GP}(n-1)}$ et
appelé le \uline{groupe projectif de degré} $\symup{n-1}$ \uline{sur} $\symup{\text{\uwave{Z}}}$. Il représente
le foncteur qui, à $\symup{S}$, associe le groupe des sections du faisceau
$\symup{\text{\uwave{Gl}}(n)_S/
\text{\uwave{Gl}}(1)_S}$, où
$\symup{\text{\uwave{Gl}}(n)_S}$ désigne le faisceau des germes de sections
de $\symup{\operatorname{Gl}(n)_S}$ sur $\symup{S}$. (On fera attention que les sections de
$\symup{\operatorname{GP}(n-1)_S}$ sur $\symup{S}$ ne proviennent pas en général de sections
de $\symup{\operatorname{Gl}(n)_S}$ sur $\symup{S}$!). Notons que l'on peut prouver que
$\symup{\operatorname{GP}(n-1)}$ représente également le foncteur
$\symup{S\rightsquigarrow\operatorname{Aut}_S(P_S^{n-1})}$, (où
$\symup{P_S^{n-1}}$ est le schéma projectif-type de dimension relative $\symup{n-1}$
sur $\symup{S}$), du moins pour $\symup{S}$ noethérien. C'est de cette façon qu'il s'introduit
en théorie des modules.

Soit $\symup{S}$ un schéma noethérien, qu'on peut si on veut supposer affine, et soit
$\symup{X}$ un $\symup{S}$-préschéma quasi-projectif, muni d'un faisceau inversible $\symup{\mathcal L}$
très ample relativement à $\symup{S}$. On suppose que le groupe
$\symup{G=\operatorname{GP}(n)_S}$ opère sur $\symup{X}$ et en même temps sur $\symup{\mathcal L}$ (de
façon compatible avec ses opérations sur $\symup{X}$), et qu'il opère \uline{librement}
sur $\symup{S}$.
\src{117}

\result{CONJECTURE}{8.1} Sous les conditions précédentes:

\smallskip
\noindent $\symup{1^\circ}$ \uline{La relation d'équivalence définie par} $\symup{G}$ \uline{est
effective, le quotient} $\symup{Y=X/G}$ \uline{est de type fini sur} $\symup{S}$ \uline{et le morphisme
canonique} $\symup{f:X\to Y}$ \uline{est plat et surjectif} (\uline{donc} $\symup{X}$ \uline{devient un fibré
principal homogène sur} $\symup{Y}$, \uline{de groupe} $\symup{G\times_S Y=\operatorname{GP}(n)_Y}$).

\smallskip
\noindent $\symup{2^\circ}$ \uline{Soit} $\symup{\mathcal L'}$ \uline{le faisceau inversible sur} $\symup{Y}$ \uline{déduit de} $\symup{\mathcal L}$ \uline{par \fquote{descente fidèlement plate} par} $\symup{f}$ (cf. [1], B, \S~th. 1). \uline{Alors} $\symup{\mathcal L'}$ \uline{est \fquote{préample} sur} $\symup{Y}$ \uline{par rapport à} $\symup{S}$,
i. e. \uline{il existe un entier} $\symup{m}$ \uline{et un morphisme quasi-fini de} $\symup{Y}$ \uline{dans un schéma projectif-type convenable} $\symup{P_S^N}$, \uline{tels que} $\symup{(\mathcal L')^{\otimes m}}$ \uline{soit isomorphe à l'image inverse de} $\symup{\mathcal O_{P_S^N}(1)}$.

On notera que même si $\symup{X}$ est séparé sur $\symup{S}$, il pourra arriver que $\symup{Y}$ ne soit
pas séparé sur $\symup{S}$ (situation qui se rencontre dans des situations de
\fquote{problèmes de modules} pas du tout pathologiques). Si $\symup{1^\circ}$ est vérifié,
alors $\symup{Y}$ est séparé si et seulement si la relation d'équivalence définie par
$\symup{G}$ a un graphe fermé, i. e. si $\symup{G\times X\to X\times X}$ a une image fermée
(c'est alors une immersion fermée). Si $\symup{Y}$ est séparé, alors $\symup{\mathcal L'}$ est
préample sur $\symup{Y}$ par rapport à $\symup{S}$ si et seulement s'il est ample, i. e. si
une puissance tensorielle convenable définit une immersion projective. Dans
les problèmes de modules mentionnés dans l'introduction, on peut montrer que la
relation d'équivalence à laquelle on parvient à bien un graphe fermé.

\medskip
\result{REMARQUES}{8.1} Nous avons supposé
$\symup{G=\operatorname{GP}(n)_S}$ pour fixer les idées et parce que c'est le cas le
plus important en pratique, à l'heure actuelle. L'hypothèse raisonnable à
faire sur $\symup{G}$ semblerait plutôt que $\symup{G}$ soit une des \fquote{formes} sur $\symup{S}$ d'un
des groupes de Tohokû (dont la construction sur les entiers a été faite
également par CHEVALLEY). Le seul fait positif qui me soit connu dans la
direction de la conjecture précédente est le suivant: \uline{Soit} $\symup{X}$ \uline{un schéma
affine sur un corps} $\symup{k}$ \uline{de caractéristique} $\symup{0}$, \uline{où le groupe}
$\symup{\operatorname{Gl}(n)_k}$ \uline{ou} $\symup{\operatorname{GP}(n-1)_k}$ \uline{opère librement. Alors
la relation d'équivalence définie par} $\symup{G}$ \uline{est effective, le quotient} $\symup{X/G}$ \uline{est
également affine, et le morphisme} $\symup{X\to X/G}$ \uline{est plat et surjectif}. La
démonstration utilise le fait suivant (qui pour l'instant n'est démontré qu'en
caractéristique $\symup{0}$): lorsque l'on fait opérer $\symup{G}$ sur l'anneau affine $\symup{A}$ de
$\symup{G}$, considéré comme espace vectoriel sur $\symup{k}$, la représentation triviale de
$\symup{G}$ n'intervient qu'une fois (dans une suite de composition d'un sous-espace
vectoriel de dimension finie sur $\symup{k}$ stable sous $\symup{G}$). Il semble possible
qu'une utilisation systématique de la
\src{118}
théorie des représentations linéaires de $\symup{G}$ finira par donner une
démonstration de la conjecture, du moins lorsque l'on est sur un corps de base.
Lorsque l'on n'est plus sur un corps de base, le conférencier ignore tout.

\begin{center}BIBLIOGRAPHIE\end{center}
\addcontentsline{toc}{section}{Bibliographie}

\begin{description}
\item[{\normalfont[1]}] GROTHENDIECK (Alexander). - Techniques de descente et théorèmes
d'existence en géométrie algébrique, I: Généralités, Séminaire Bourbaki,
1959/60, n\textsuperscript{o} 190, 29 p.

\item[{\normalfont[2]}] GROTHENDIECK (Alexander). - Techniques de descente et théorèmes
d'existence en géométrie algébrique, II: Le théorème d'existence en théorie
formelle des modules, Séminaire Bourbaki, 1959/60,
n\textsuperscript{o} 195, 22 p.

\item[{\normalfont[3]}] GROTHENDIECK (A.) et DIEUDONNÉ (J.). - Éléments de géométrie
algébrique, I: Le langage des schémas. - Paris, Presses universitaires de
France, 1960 (Institut des hautes Études scientifiques, Publications
mathématiques, 4).
\end{description}

\begin{center}
  \rule{.14\linewidth}{.3pt}
\end{center}
